% 10_body.tex — Day2 산출물 5/5 (⑩) 리스크 등록부(정량화)와 허용범위 격자 (빈 템플릿 / 완성 예시 공용 본문)
% 드라이버: 10_리스크등록부_빈템플릿.tex (\wbblanktrue) / 10_리스크등록부_완성예시.tex (\wbblankfalse)
% 내용 출처: PM트랙/Day2_워크북.md §5.3(항목 가이드) §5.4(빈 템플릿) §5.5(온담 P1 완성 예시본)
\documentclass[10pt,a4paper]{article}
\usepackage[a4paper,margin=16mm,top=14mm,bottom=20mm]{geometry}
\def\DocNum{5}
\def\DocTitle{리스크 등록부(정량화)와 허용범위 격자}
\def\DocSub{⑩ Quantified Risk Register \& Tolerance Grid}
\def\DocVariant{\B{빈 템플릿}{온담 P1 완성 예시}}
\usepackage{wbstyle}
\begin{document}

\docheader
 {\Bg{[정식 명칭과 코드]}{차세대 주문·재고 통합 플랫폼 구축 (ONE ONDAM P1)}}
 {\Bg{[v1.x / 판정 시점]}{v1.0 / 판정 2026-03-27 (S6 종료), G1 제출}}
 {\Bg{[PM 실명 — RACI 26행 A]}{P1 PM (RACI 26행 A, 게이트 갱신)\rd{7mm}}}
 {\Bg{[스폰서 실명]}{정미란 CFO (소유자 각자 서명)\rd{7mm}}}

% ------------------------------------------------------------------ 1
\secbar{1. 문서 정보와 척도 정의}
\guide{등록부 버전, 판정 시점, 확률 척도(\%)와 영향 척도(원가 억 원 / 일정 영업일 / 편익 억 원·년)의 정의, EMV 산식, 일정 1일의 금액 환산, 관련 문서. \textbf{확률 1~5점 × 영향 1~5점 = “점수”는 순위 외에 아무 정보가 없고 합계도 컨틴전시와 비교할 수 없다.} 원가 영향과 일정 영향은 다른 것 — 두 축을 따로 잰다.}
{\footnotesize
\begin{tabularx}{\linewidth}{|L{32mm}|Y|}\hline
\hd{항목} & \hd{내용} \\ \hline
\B{\rd{9mm}}{}버전 / 판정 시점 / 입력 & \Bg{[v1.x / 날짜 / 초기 리스크 목록 vN, 원가 기준선, 일정 기준선, 격자 초안]}{v1.0 / 2026-03-27 / 초기 리스크 목록 v1.0(2026-01-09), 원가 기준선 v0.9, 일정 기준선 v0.9, 거버넌스 §4 격자 초안} \\ \hline
\B{\rd{9mm}}{}확률 척도 & \Bg{[\% 직접 기입 또는 구간 중앙값 10/25/40/60/80 — 근거 우선순위]}{\% 직접 기입. 근거 우선순위: 트리거 근접성 > 과거 이력(2019·2022) > 프리모템 득표. 구간 참고값 10 / 25 / 40 / 60 / 80} \\ \hline
\B{\rd{11mm}}{}영향 척도 & \Bg{[원가 = 예비비에서 실제 나갈 금액(억) / 일정 = 경로 지연 영업일 / 편익 = 연간 편익 감소(억/년, B-ID)]}{원가 = 컨틴전시 또는 관리예비비에서 실제로 나갈 금액(억). 일정 = 해당 경로의 지연 영업일(임계경로면 G3·G4 지연). 편익 = 연간 편익 감소(억/년, B-ID). 출처 없는 값은 [추가 설정]} \\ \hline
\B{\rd{9mm}}{}EMV / 일정 1일 환산 & \Bg{[EMV = P × 영향 / 1영업일 ≈ [\ ]억 (스프린트 비용 ÷ 영업일) — 원가 EMV에 합산하지 않음]}{EMV = P × 영향. 일정 1영업일 ≈ 0.14억(스프린트 1.41억 ÷ 10일) — 원가 EMV에 합산하지 않고 별도 관리} \\ \hline
\B{\rd{9mm}}{}작성 / 승인 & \Bg{[PM, 소유자 서명 / 스폰서]}{P1 PM, 소유자 각자 서명 / 정미란 (RACI 26행 A = P1 PM, 게이트 갱신)} \\ \hline
\end{tabularx}}

% ------------------------------------------------------------------ 2~4 (가로)
\landscapeon
\secbar{2~4. 등록부 — 정량화 · 대응 · 2차 리스크와 잔여 노출}
\guide{Day1 초기 목록의 원인·사건·영향·소유자·트리거를 그대로 잇고 확률(\%)·원가 영향(억)·일정 영향(영업일)·편익 영향·EMV를 붙인다(기회는 음수). 영향 금액의 출처를 적는다(롤백 2.1억 = 헌장 §7). 전략(회피·전가·완화·수용 / 활용·공유·강화·수용)만 적지 말고 \textbf{조치를 WBS 코드로 연결}해 기준선에 포함되어 있는지 밝힌다 — 기준선에 없는 조치는 조건부 대응이며 컨틴전시 트리거가 된다. 트리거는 관측 가능한 신호. \textbf{잔여를 0으로 적지 말 것}(대응이 리스크를 없애지 않는다). 자금원은 원인·결정 권한이 어느 계층인가로 정한다. PM이 소유자인 리스크는 3개를 넘지 않는다. 확률 100\%(이미 일어난 것 = 이슈)나 0\%는 리스크가 아니다.}
{\scriptsize\setlength{\tabcolsep}{2pt}
\begin{longtable}{|L{9mm}|L{33mm}|C{7mm}|L{12mm}|L{11mm}|L{12mm}|C{9mm}|C{9mm}|L{11mm}|L{32mm}|L{17mm}|L{20mm}|L{18mm}|L{15mm}|L{13mm}|L{10mm}|}\hline
\hd{ID} & \hd{사건 요약 (원인→사건→영향)} & \hd{P\%} & \hd{원가 영향(억)} & \hd{일정 영향(일)} & \hd{편익 영향} & \hd{원가 EMV} & \hd{일정 EMV} & \hd{전략} & \hd{조치 (WBS·기준선 포함·비용)} & \hd{소유자 (조치자)} & \hd{트리거 (신호·시점)} & \hd{2차 리스크} & \hd{잔여 P × 영향 = EMV} & \hd{자금원} & \hd{검토} \\ \hline \endhead
\ifwbblank
\rd{5mm}R-01 & \g{[원인 → 사건 → 영향]} & \g{[\%]} & \g{[억, 출처]} & \g{[일, 임계?]} & \g{[억/년, B-ID]} & \g{[P×원가]} & \g{[P×일]} & \g{[회피/전가/완화/수용]} & \g{[1.n.m, 기준선 포함 여부, 비용]} & \g{[실명 (실명)]} & \g{[관측 가능 신호, 시점]} & \g{[R-nn 원인→사건]} & \g{[P × 영향 = EMV]} & \g{[컨틴전시/관리예비비/밖]} & \g{[게이트]} \\ \hline
\rd{5mm}R-02 & & & & & & & & & & & & & & & \\ \hline
\rd{5mm}R-03 & & & & & & & & & & & & & & & \\ \hline
\rd{5mm}R-04 & & & & & & & & & & & & & & & \\ \hline
\rd{5mm}R-05 & & & & & & & & & & & & & & & \\ \hline
\rd{5mm}R-06 & & & & & & & & & & & & & & & \\ \hline
\rd{5mm}R-07 & & & & & & & & & & & & & & & \\ \hline
\rd{5mm}R-08 & & & & & & & & & & & & & & & \\ \hline
\rd{5mm}R-09 & & & & & & & & & & & & & & & \\ \hline
\rd{5mm}R-10 & & & & & & & & & & & & & & & \\ \hline
\rd{5mm}R-11 & & & & & & & & & & & & & & & \\ \hline
\rd{5mm}R-12 & & & & & & & & & & & & & & & \\ \hline
\rd{5mm}R-13 & & & & & & & & & & & & & & & \\ \hline
\rd{5mm}\textbf{O-01} & \g{(기회)} & & \g{(−)} & \g{(−)} & & \g{(−)} & \g{(−)} & \g{[활용/공유/강화/수용]} & & & & & \g{(−)} & \g{—} & \\ \hline
\rd{5mm}\textbf{O-02} & \g{(기회)} & & \g{(−)} & \g{(−)} & & \g{(−)} & \g{(−)} & & & & & & \g{(−)} & \g{—} & \\ \hline
\else
R-01 & P2 리드타임 41주 → MC 2027-04 → 물류 인터페이스 시험 G3 후 → 컷오버 물류 분리 또는 창 이탈 & 60 & 4.5 (2차 컷오버·재시험 [추가 설정]) & 20 (임계, A15) & — & 2.70 & 12.0 & 완화 & 분리 가능 설계 1.3.4·1.6.9(기준선 포함, 추가 FP 0), 분리 결정 G1 이사회 부의(RACI 24행) & 프로그램 PMO장 (P1 PM) & T0 2026-04-24 리드타임 확정 통보; 2026-10 P2 진도 & R-14 두 번 컷오버 → 매장 혼선·2차 창 성수기 접근 & 60 × 2.0 = \textbf{1.20} & 관리예비비 & G1, 2026-10 \\ \hline
R-02 & 가맹 388점 동의 지연·거부 → SW 배포 범위 축소 → 가맹 데이터 미수집 & 50 & 3.0 (분리 배포·재작업) & 15 (비임계, 1.1.6) & B-02·B-04 가맹분 약 −12/년 & 1.50 & 7.5 & 완화 & 포털·POS SW 분리 릴리스(O-02, 1.4.3 기준선 포함), 데이터 용도 명세서·단말 조사 1.1.6(기준선 포함) & 오정근 (매장운영팀장) & 협의회 서면 질의; 2026-07 동의율 중간 집계 60\% 미만 & R-15 포털 선행 릴리스 결함 → 협상 지렛대 상실 & 30 × 2.0 = \textbf{0.60} & 컨틴전시 0.60 & 2026-06, G1 \\ \hline
R-03 & 대성 준실시간 거부 → 센터 재고 야간 배치 → B-02 매장 한정 & 40 & 2.0 (폴백 확정·재시험) & 10 (비임계, A08 TF 20) & B-02 센터분 미산정 & 0.80 & 4.0 & 완화+수용 & 폴백 모듈 1.3.3(기준선 포함), 부하 산정서 2026-03(기준선 포함) & 한동석 (센터운영 담당) & 협상 착수 2026-03; 대성 서면 회신; 06-30 미타결 & R-16 폴백 2회 배치 부하 → 야간 작업 창 충돌 & 40 × 1.2 = \textbf{0.48} & 컨틴전시 0.50 & G1 \\ \hline
R-04 & CFO 148억 확정 → P1 2026 배분 축소 → 라이선스·선약정 지연·인력 축소 → 창 이탈 & 60 & 5.0 (종량제 프리미엄·대기·압축 [추가 설정]) & 20 (임계, A03→A05) & — & 3.00 & 12.0 & 완화 & 이월 vs 총액 삭감 시나리오별 축소안 2026-01-30 제출(완료), PV 봉우리 이월 옵션(⑨ 항목 5) & 정미란 (P1 PM·재경팀장) & 재무본부 자금계획 2026 확정(2026-02); 배분 통보 & R-17 이월분 2027 집중 → 2027 자금 계획 초과 & 40 × 2.5 = \textbf{1.00} & 관리예비비 & 2026-02, G1 \\ \hline
R-05 & 겸임 인력 참석 저하 → 미결정 대기 → 변경대가 청구 & 70 & 2.4 (대기 비용, 스프린트 1.7회분) & 10 (임계, A01·A05) & — & 1.68 & 7.0 & 완화 & 결정 응답 SLA 5영업일(RACI·거버넌스 §3 병행 행), 대체 결정권자 지정, 투입률 월 측정(1.6.2) & 박세연 (P1 PMO) & 회의 참석률 70\% 미만 2주 연속; 변경협의체 대기 비용 상정 & — & 40 × 1.5 = \textbf{0.60} & 컨틴전시 0.60 & 월, G1 \\ \hline
R-06 & OD-POS 개발자 1명 이탈 → 13본 분석 불가 → 전환 매핑·G2 지연 & 30 & 3.5 (리버스 엔지니어링 외주) & 25 (근접 임계, A02→A09 TF 5) & — & 1.05 & 7.5 & 완화 & 13본 복원 S2–S6 전진 배치 1.5.5(기준선 포함, O-01), 지식 이전 세션 기록, 인사본부 유지 협의 & 박세연 (IT본부 운영파트장) & 퇴직 의사 표명; 2026-03 문서화 진도 50\% 미만 & — & 20 × 1.5 = \textbf{0.30} & 컨틴전시 0.30 & G1 \\ \hline
R-07 & 동결 직전 요구 집중 → ±35건 초과 → FP 재산정·기준선 재확인 & 60 & 1.5 (재산정 공수·설계 지연) & 10 (임계, 구현 스프린트) & — & 0.90 & 6.0 & 완화 & 접수 마감 동결 2주 전(일정 기준선 캘린더 포함), Must ≤ 60\%(RTM 규칙), 오정근 Should 조정 & P1 PM (PMO 요구사항 담당) & 월 접수 건수 추이; 순변동 ±35건 초과 & — & 40 × 1.0 = \textbf{0.40} & 컨틴전시 0.40 & 매월 \\ \hline
R-08 & 앱 벤더 계약 외 주장·별도 견적 4억 → 앱 주문 연동 지연 → B-05·B-01 지연 & 50 & 4.0 (프리모템 원문 8; 벤더 견적은 앱 유지보수 계약, P1측 재작업 포함) & 15 (비임계, 1.2.3) & B-05 37/년 M+24 이후로 & 2.00 & 7.5 & 전가+완화 & 앱 유지보수 계약에 소스 인수·API 연동 조항 추가 2026-03(RACI 12행), API 규격 조기 제공 2026-04(1.2.3 기준선 포함) & 박세연 (계약: 오정근) & 벤더 API 규격 회신 4주 지연; 계약 변경 미서명 2026-04 & R-18 계약 변경 협상이 유지보수 단가 인상 요구로 번짐 & 30 × 2.0 = \textbf{0.60} & 컨틴전시 0.60 (벤더 견적은 P1 예산 외) & G1 \\ \hline
R-09 & 별도 법인 개인정보 처리 주체 미확정 → 공장 연동 규제 범위 제외 또는 P1 흡수 & 40 & 1.8 (공장 연동 재설계) & 10 (비임계, 1.3.5) & — & 0.72 & 4.0 & 완화 & 별도 법인 위탁 계약 초안 2026-03(RACI 13행), P4 규제 아키텍처 연동 체크포인트(1.5.1) & 최영주 (P1 PM, 조성태) & P4 위탁 구조 2026-05 미확정 & — & 20 × 1.5 = \textbf{0.30} & 컨틴전시 0.30 & G1 \\ \hline
R-10 & G2 정본 판정 후 마스터 코드 변경 → 610 테이블 재매핑 3.8억·6주 → 리허설 지연 & 25 & 3.8 (헌장 §7) & 30 (근접 임계, A13 TF 5) & — & 0.95 & 7.5 & 회피 & 마스터 변경 동결 규칙 G2 전 CCB 의결(1.5.8 산출물 ④), 판정 후 변경은 R5 이후 릴리스로 & 박세연 (데이터 오너 4인) & G2 전 마스터 변경 요청 건수; 신규 브랜드·B2B 코드 요구 & — & 10 × 3.8 = \textbf{0.38} & 컨틴전시 0.40 (전액 실현 시 CFO 병행) & G2 \\ \hline
R-11 & SAC 미합의 → G4 인수 거부 → 하자보수 체제 장기화 & 40 & 2.5 (PMO·수행사 잔류) & 40 (임계, A20) & — & 1.00 & 16.0 & 완화 & SAC 초안 G1 공동 작성(RACI 22행), G3 SAC 사전 점검, 운영 기술 문서 1.6.10(기준선 포함) & 박준영 (P1 PM) & G1 SAC 초안 미합의; G3 사전 점검 미해소 항목 20건 초과 & — & 20 × 2.0 = \textbf{0.40} & 컨틴전시 0.40 & G1, G3 \\ \hline
R-12 & 편익 분모 분쟁·노조 실사 비협조 → 편익 판정 불가 & 40 & 1.0 (재실사·표본) & 10 (비임계) & 편익 판정 자체 & 0.40 & 4.0 & 완화 & 편익 책임자 서명 G1(BC §9), 두 분모 리포트 StRS-034(기준선 포함), “P1 편익에 인건비 없음” 문서화 → 인사본부 채널 & 오정근·한동석 (나영선 협의) & M+3(2026-04) 분모 미확정; 실사 협조 거부 & — & 25 × 0.8 = \textbf{0.20} & 컨틴전시 0.20 & G1 \\ \hline
R-13 & 구형 단말 펌웨어 실패 → 일부 매장 수기 복귀·롤백 & 35 & 2.1 (롤백 2.1억, 헌장 §7) & 6 (롤백 6영업일) & — & 0.74 & 2.1 & 완화 & 단말 전수 조사 2026-04·비호환 대응안 1.1.6(기준선 포함), 파일럿 구형 기종 포함 1.1.7 & P1 PM (오정근 조사) & 조사 결과 비호환 10\% 초과(2026-04-30) & — & 15 × 2.1 = \textbf{0.32} & 컨틴전시 0.30 & G1, 2027-01 \\ \hline
\textbf{O-01} & 13본 문서화 조기 완료 → 전환 매핑 가속 → G2 조기·리허설 여유 & 50 & −0.8 (리허설 추가 회차 불필요) & −5 & — & −0.40 & −2.5 & 활용 & R-06 조치와 통합(1.5.5 S2–S6) & 박세연 & 문서화 진도 2026-03 80\% 이상 & — & 50 × −0.8 = \textbf{−0.40} & — & G1 \\ \hline
\textbf{O-02} & 포털 선행 릴리스(2026-Q4) → 부속서 협상 지렛대 → R-02 확률 감소·B-08 조기 & 60 & −1.0 (R-02 대응 비용 절감) & −5 & B-08 M+12 → M+9 & −0.60 & −3.0 & 활용 & 1.4.3 프로토타입 2026-05·선행 릴리스 S23–S24(기준선 포함), 헌장 §13에서 포털 이연 불가 & 오정근 (P1 PM) & 프로토타입 시연 2026-05 협의회 반응 & R-15 (R-02와 공유) & 60 × −1.0 = \textbf{−0.60} & — & 2026-05, G1 \\ \hline
\fi
\end{longtable}}
\B{\small\g{2차 리스크는 예비 번호(R-nn)로 등록하고 트리거 발동 시 정식 정량화한다. 위협 EMV 합 \hspace{12mm} / 기회 EMV 합 \hspace{12mm} / 잔여 위협 EMV 합 \hspace{12mm}}\par}{\small 2차 리스크 R-14~R-18은 예비 번호로 등록하고 트리거 발동 시 정식 정량화한다. R-14(두 번 컷오버)는 R-01 분리 결정 시 즉시 정량화 대상. 위협 13건: 회피 1(R-10), 전가+완화 1(R-08), 완화 10, 완화+수용 1(R-03). 기준선 포함 조치 11건, 조건부 대응 6건. 소유자 PM 2건(R-07, R-13).\par}

% ------------------------------------------------------------------ 5 (가로)
\clearpage
\secbar{5. EMV 합계와 컨틴전시·허용범위 대조}
\guide{대응 전 합, 대응 후(잔여) 합, 컨틴전시 배정(⑨ 항목 6)과의 대조, 프로젝트 계층 리스크 허용범위와의 대조. 확인식: “잔여 컨틴전시 대상 EMV ≤ 컨틴전시 배정 ≤ 컨틴전시”, “잔여 위협 EMV 합 ≤ 허용범위”. 넘으면 (a) 대응 강화 (b) 관리예비비 이관 (c) 허용범위 재협상 중 무엇을 택했는지 적는다 — \textbf{맞추기 위해 확률을 낮추지 말 것.} EMV는 평균이며 꼬리를 덮지 못한다 — 그 한계를 한 문단으로.}
{\footnotesize
\begin{tabularx}{\linewidth}{|L{30mm}|C{22mm}|C{20mm}|C{18mm}|L{74mm}|Y|}\hline
\hd{구분} & \hd{위협 합} & \hd{기회 합} & \hd{순} & \hd{비교 대상} & \hd{판정} \\ \hline
\ifwbblank
\rd{9mm}대응 전 & & & & \g{리스크 허용범위 [\ \ ]억 / 컨틴전시 [\ \ ]} & \\ \hline
\rd{9mm}대응 후 (잔여) & & & & \g{허용범위 [\ \ ]억} & \g{[이내/초과 — 여유]} \\ \hline
\rd{11mm}자금원별 잔여 & \multicolumn{3}{l|}{\g{컨틴전시 [\ \ ] (n건) / 관리예비비 [\ \ ] / 기준선 밖 [\ \ ]}} & \g{컨틴전시 배정 [\ \ ] ≤ 컨틴전시 [\ \ ] / 관리예비비 [\ \ ]} & \\ \hline
\rd{9mm}대응의 가치 & \g{[전 − 후]} & & & \g{대응 조치 비용: 기준선 포함 n건 / 조건부 n건} & \\ \hline
\else
대응 전 & \textbf{17.44} & −1.00 & 16.44 & 리스크 허용범위 8억 / 컨틴전시 10.2 & 초과 — 대응 없이는 예외 보고 대상 \\ \hline
대응 후 (잔여) & \textbf{6.78} & −1.00 & 5.78 & 허용범위 8억 & 이내 (여유 1.22 — 단일 리스크 하나 실현되면 소진) \\ \hline
자금원별 잔여 & \multicolumn{3}{l|}{컨틴전시 대상 4.58 (11건) / 관리예비비 대상 2.20 (R-01·R-04) / 기준선 밖 0} & 컨틴전시 배정 4.60 ≤ 컨틴전시 10.20 / 관리예비비 12.0 & 충족 \\ \hline
대응의 가치 & 17.44 − 6.78 = 10.66 감소 & & & 대응 조치 비용: 기준선 포함 11건(추가 원가 0), 조건부 6건(컨틴전시) & 대응은 기준선 안의 작업패키지로 대부분 흡수 \\ \hline
\fi
\end{tabularx}}
\vspace{2mm}\noindent\textbf{\small EMV의 한계} \g{(한 문단 — 리스크 간 의존·전파 경로, 꼬리 분포, 처방)}\par\vspace{1mm}
\begin{fieldbox}
\ifwbblank\lines{4}{8mm}\else
6.78은 평균이다. 13개 위협이 독립이 아니고(R-04 삭감 → R-05 인력 → R-07 요구 집중, R-01 P2 → R-14 두 번 컷오버 → R-13 단말), 결합 시 영향은 합보다 크다. IT 원가 초과의 분포는 꼬리가 두꺼워(Flyvbjerg 2022) 컨틴전시를 늘려도 꼬리를 덮지 못한다. 그래서 이 등록부의 처방은 컨틴전시 증액이 아니라 전파 경로 절단이다 — 컷오버 매장·주문/물류 분리(1.6.9), 릴리스 R5 이연(헌장 §13), 마스터 동결(R-10 회피). 이 세 조치의 공통점은 한 사건이 전체를 밀지 못하게 하는 것이다. G2·G3에서 몬테카를로(3점 추정 기반) 실시 여부는 프로그램 PMO 판단 [추가 설정].
\fi
\end{fieldbox}

% ------------------------------------------------------------------ 6 (가로)
\secbar[60mm]{6. 일정 리스크 노출}
\guide{일정 영향 EMV 합(대응 전·후)과 일정 허용범위·버퍼(합류 버퍼, 규제 버퍼)와의 대조. 임계경로에 걸리는 리스크와 걸리지 않는 리스크를 구분한다 — \textbf{병렬 경로의 지연은 더해지지 않는다}(비임계 경로의 지연은 TF 안에서 흡수). ⑧ 항목 5의 TF 열을 참조. 원가만 정량화하고 일정은 “지연 가능”으로 두지 말 것.}
{\footnotesize
\begin{tabularx}{\linewidth}{|L{34mm}|C{26mm}|L{70mm}|L{50mm}|Y|}\hline
\hd{구분} & \hd{일정 EMV 단순 합} & \hd{임계경로 직접 영향 (해당 리스크 ID와 합)} & \hd{비교 대상 (허용범위 + 버퍼)} & \hd{판정} \\ \hline
\ifwbblank
\rd{10mm}대응 전 & \g{[영업일]} & \g{[R-nn … = 합]} & \g{허용범위 +[\ ] + 합류 버퍼 [\ ] = [\ ]} & \\ \hline
\rd{10mm}대응 후 (잔여) & & & & \g{[여유 n일 — 단독으로 창을 깨는 리스크?]} \\ \hline
\else
대응 전 & 97.1영업일 & R-01·R-04·R-05·R-10·R-11: 12.0 + 12.0 + 7.0 + 7.5 + 16.0 = 54.5 & 허용범위 +15 + G3 합류 버퍼 15 = 30 & 초과 — 대응 필수 \\ \hline
대응 후 (잔여, [추가 설정] 잔여 일수) & 33.0영업일 & 6.0 + 4.0 + 2.0 + 3.0 + 4.0 = 19.0 & 30 & 이내, 여유 11일. 단 R-01 실현 시 단독 20일 → 컷오버 창 이탈 \\ \hline
\fi
\end{tabularx}}
\B{\small\g{일정 노출 / 원가 노출 비교와 조정 변수(범위)의 필요성: \hrulefill}\par}{\small 일정 노출(19 / 30 = 63\%)이 원가 노출(6.78 / 8 = 85\%)보다 여유가 있어 보이지만, 컷오버 창은 허용범위 0이므로 R-01 하나가 창을 깬다. 헌장 §13이 범위를 조정 변수로 둔 이유가 여기서 숫자로 확인된다 — 일정을 지키려면 범위(R5 이연)를 쓴다.\par}

% ------------------------------------------------------------------ 7 (가로)
\clearpage
\secbar{7. 허용범위 격자 확정본}
\guide{7축(편익·원가·일정·품질·범위·지속가능성·리스크) × 3계층(작업패키지·프로젝트·프로그램)의 값과 \textbf{각 축의 측정 정의}(무엇을 언제 어떻게 재는가 — 분모·시점·집계), 초과 시 보고 경로, 사내 규정과의 관계, Day1 초안 대비 변경 표시. 값 뒤에 괄호로 분모·시점을 적는다. 마일스톤별 예외(컷오버 창 0)를 일정 축에. 숫자만 있고 정의가 없으면(“0.4건/FP”가 어느 시점 어떤 결함인지) 판정할 수 없다. 고정가 계약에서 수행사 원가는 발주사가 관리할 대상이 아니다 — 작업패키지 계층이 계약과 맞는지 확인.}
{\scriptsize
\begin{longtable}{|L{13mm}|L{34mm}|L{47mm}|L{27mm}|L{48mm}|L{24mm}|L{26mm}|L{26mm}|}\hline
\hd{축} & \hd{작업패키지 (수행사 PM / 패키지 담당)} & \hd{프로젝트 (PM)} & \hd{프로그램 (스폰서)} & \hd{측정 정의 (분모·시점·집계)} & \hd{초과 예상 시 보고} & \hd{사내 규정} & \hd{변경 (Day1 초안 대비)} \\ \hline \endhead
\ifwbblank
\rd{13mm}편익 & \g{—} & \g{[−n\%p]} & \g{[−n\%p]} & \g{[실현률 전망, 기준 편익액, 시점, 책임자 / 중단 조건과 구분]} & \g{[경로·시한]} & & \\ \hline
\rd{13mm}원가 & \g{[발주 집행 패키지 +n\% / SI 패키지 FP ±n\%]} & \g{[+n억 = 계획 원가의 +n\% — 분모 명시]} & \g{[+관리예비비]} & \g{[EAC 기준 예상 초과]} & & \g{[단일 건 기준 금액, 내부통제]} & \\ \hline
\rd{13mm}일정 & \g{[+n영업일]} & \g{[종료 +n / 고정 마일스톤 0 / 중간 게이트 +n 보고]} & \g{[+n (버퍼 포함)]} & \g{[CPM 완료일 대비 예상, TF 기준]} & & & \\ \hline
\rd{13mm}품질 & \g{[결함밀도]} & \g{[n건/FP = 정의 + 심각도 미해결 0]} & & \g{[시험 결과서, 결함 도구, 감리]} & & & \\ \hline
\rd{13mm}범위 & \g{[±n\% FP]} & \g{[±n\% = ±n건 및 FP ±n점]} & \g{[±n\%]} & \g{[월 동결 순변동, FP 누적]} & & \g{[계약 부속서]} & \\ \hline
\rd{13mm}지속가능성 & & \g{[0 = 무엇의 0인가]} & & & \g{[즉시]} & & \\ \hline
\rd{13mm}리스크 & \g{[잔여 EMV n억]} & \g{[잔여 위협 EMV 합 n억 + 단일 n억 초과 보고]} & \g{[n억]} & \g{[월 갱신 시 잔여 EMV 합]} & & & \\ \hline
\else
편익 & — & −5\%p & −10\%p & G1·G2·G4 시점 편익 실현률 \textbf{전망}(직접 편익 118억 기준, BC §9 책임자 산정). 손익분기 66\% 미만 전망은 허용범위가 아니라 중단 조건(헌장 §14) & 스폰서 5영업일 & 판정 권한 재무본부 & 측정 시점 추가 \\ \hline
원가 & 발주 집행 패키지(조달·단말·전환·인프라): WBS 사전 추정 +2\% / SI 패키지: FP ±1\% & \textbf{+10.2억 = 컨틴전시 전액 = 계획 원가 145.8의 +7.0\%} (계획 원가 + 허용범위 = 집행 한도 156.0) & +12.0억 = 관리예비비 (= 승인 예산 168.0) & EAC(Day3 EVM) 기준 집행 한도 초과 \textbf{예상} 시 & 스폰서 5영업일, 결정 10영업일 & 5억 초과 단일 건 CFO 병행, 5,000만 원 이상 내부통제 지침 & \textbf{초안 +6\%(10.08, 승인 예산 기준) → 컨틴전시 전액} \\ \hline
일정 & +5영업일 (패키지 완료일) & G4 기준 \textbf{+15영업일} / \textbf{컷오버 창·규제 기한 0} / G1·G2·리허설 3회차 +5(보고) & +30영업일 (G3 합류 버퍼 15 포함) & 일정 기준선 CPM 완료일 대비 예상 완료일, 임계경로 TF 기준 & 스폰서 5영업일; 창 이탈 예상은 즉시 & — & \textbf{마일스톤별 예외 추가} \\ \hline
품질 & 결함밀도 0.13건/FP (패키지 통합시험) & \textbf{0.4건/FP} = 통합시험~UAT 발견 결함 총수 ÷ 12,400 + \textbf{심각도 1·2 미해결 0건(G3)} & 0.8건/FP & 시험 결과서·결함 관리 도구, 감리 확인 & 스폰서 (G3 전) & 감리 시정조치는 허용범위 무관 이행 & \textbf{측정 정의 추가} \\ \hline
범위 & ±1\% (패키지 FP) & \textbf{±3\% = ±35건(1,180 기준) 및 FP ±372점} & ±6\% & 월 동결 시 순변동(건), FP 누적 변동(계약 재확인 트리거와 동일) & 변경협의체 합의 → CCB & 계약 부속서 1 & \textbf{FP 기준 병기} \\ \hline
지속가능성 & 0 & \textbf{0 = 개인정보 유출 사고 0, 보안진단 High 미조치 0(G3), 안전·환경 사고 0} & 0 & 보안진단 보고, 사고 보고 & 즉시 (최영주·스폰서) & 규제 보고 의무 & \textbf{정의 추가} \\ \hline
리스크 & 잔여 EMV 2.7억 (패키지 합) & \textbf{잔여 위협 EMV 합 8억} (현재 6.78) \textbf{+ 단일 리스크 잔여 EMV 2억 초과 시 스폰서 보고} & 16억 & 등록부 월 갱신 시 잔여 EMV 합계(컨틴전시·관리예비비 대상 합산) & 스폰서 5영업일 & 이사회 부의 금액 & \textbf{“잔여” 기준 명시, 단일 건 기준 추가} \\ \hline
\fi
\end{longtable}}

% ------------------------------------------------------------------ 8 (가로)
\clearpage
\secbar{8. Day1 초안 대비 변경과 이유}
\guide{격자에서 바뀐 칸, 초안 값, 확정 값, 이유, 수정 대상 문서·항목(헌장 §12, 거버넌스 §4 …). \textbf{이유는 오늘의 산출물(원가 기준선의 3층, 일정 기준선의 제약, RTM의 FP 기준)에서 나온 것}이어야 한다. 변경이 없는 칸은 적지 않는다. 초안을 그대로 확정했다면 오늘 계산한 것이 초안과 맞는지 확인하지 않았다는 뜻이다. 이 표가 Day1 문서로 되돌아가는 헌장 v1.1 수정 요청의 재료다.}
{\footnotesize
\begin{longtable}{|L{26mm}|L{44mm}|L{50mm}|L{92mm}|L{40mm}|}\hline
\hd{축·칸} & \hd{초안 (Day1 거버넌스 §4)} & \hd{확정} & \hd{이유 (오늘의 어느 산출물에서)} & \hd{수정 대상 문서·항목} \\ \hline \endhead
\ifwbblank
\rd{12mm}\g{[축·계층]} & \g{[초안 값]} & \g{[확정 값]} & \g{[⑥~⑩의 어느 표·어느 숫자가 초안과 어긋났는가]} & \g{[헌장 §n, 거버넌스 §n]} \\ \hline
\rd{12mm} & & & & \\ \hline
\rd{12mm} & & & & \\ \hline
\rd{12mm} & & & & \\ \hline
\rd{12mm} & & & & \\ \hline
\rd{12mm} & & & & \\ \hline
\rd{12mm} & & & & \\ \hline
\rd{12mm} & & & & \\ \hline
\else
원가·프로젝트 & +6\% (10.08억, 승인 예산 168 기준) & +10.2억 = 컨틴전시 전액 (계획 원가의 +7.0\%) & 원가 기준선 3층(⑨)에서 계획 원가 + 10.08 = 155.88 ≠ 집행 한도 156.0. 허용범위 = PM이 보고 없이 쓰는 돈 = 컨틴전시로 정의해야 층이 맞물림 & 헌장 §12 표, 거버넌스 §4, 프로그램 관리 규정 별표 해석 \\ \hline
원가·작업패키지 & +2\% & 발주 집행 패키지 +2\% / SI 패키지 FP ±1\% & 고정가 계약에서 수행사 원가는 발주사 관리 대상이 아님. 발주사가 관리하는 것은 FP와 발주사 집행 패키지 & 거버넌스 §3 1단계 조건 \\ \hline
일정·프로젝트 & +15영업일 & +15영업일 + 컷오버 창·규제 기한 0 + 중간 게이트 +5 보고 & 일정 기준선(⑧)의 제약 일자 — 창 이탈은 15일이 아니라 6개월(다음 창)이므로 허용범위 개념이 성립하지 않음 & 헌장 §12, §9 제약 \\ \hline
범위·프로젝트 & ±3\% (±35건) & ±35건 및 FP ±372점 병기 & RTM(⑦) Must 비율이 FP 기준이고 계약 재확인 트리거가 FP 누적 ±3\%. 건수만으로는 계약과 어긋남 & 헌장 §12, 범위 기술서 §6 \\ \hline
품질·프로젝트 & 0.4건/FP & 정의: 통합시험~UAT 발견 결함 ÷ FP + 심각도 1·2 미해소 0건 & 정의 없는 0.4는 판정 불가 — 미해결 기준이면 4,960건이 되어 무의미. 발견 밀도 + 미해결 0건 두 기준으로 분리 & 헌장 §3 성공 기준 행 \\ \hline
리스크·프로젝트 & EMV 8억 & 잔여 위협 EMV 합 8억 + 단일 2억 초과 보고 & 대응 전 17.44는 어느 프로젝트나 넘는다. 잔여 기준이어야 대응의 가치가 보이고, 단일 건 기준이 없으면 R-04(3.00) 하나가 합계 안에 숨는다 & 헌장 §12, 거버넌스 §4 \\ \hline
지속가능성 & 0 (안전·환경 사고) & 개인정보 유출 0, High 취약점 미조치 0, 안전·환경 0 & SI 프로젝트에서 “안전·환경”만으로는 측정 대상이 없음. 규제 축(HR-06)과 연결 & 거버넌스 §4 \\ \hline
편익 & −5\%p & 값 유지, 측정 시점(G1·G2·G4 전망) 추가, 중단 조건(66\%)과 구분 & BC §6 손익분기 66\%와 허용범위 −5\%p의 관계가 초안에 없었음 & 헌장 §12·§14 \\ \hline
\fi
\end{longtable}}
\landscapeoff

% ------------------------------------------------------------------ 9
\secbar[70mm]{9. 검토 주기와 에스컬레이션}
\guide{등록부 갱신 주기(월·게이트), 트리거 관측 담당, 리스크 허용범위 초과 시 보고 경로(헌장 §12·거버넌스 §3 인용), 신규 리스크 등록 절차(월 동결 시), 리스크 실현(이슈 전환) 시 자금 집행 순서, 소유자 서명. 게이트마다 등록부 버전이 올라간다. G1에 만들고 G4까지 안 고치는 등록부, 아무도 보지 않는 트리거를 막는 것이 이 항목이다.}
\begin{fieldbox}
\ifwbblank
\g{갱신 주기: \hrulefill}\par\vspace{4mm}
\g{트리거 관측 담당: \hrulefill}\par\vspace{4mm}
\g{초과 예상 시 보고 (경로·시한·형식): \hrulefill}\par\vspace{4mm}
\g{관리예비비 대상 리스크 트리거 발동 시: \hrulefill}\par\vspace{4mm}
\g{신규 리스크 등록 절차: \hrulefill}\par\vspace{4mm}
\g{실현(이슈 전환) 시 자금 집행 순서: \hrulefill}\par\vspace{4mm}
\g{소유자 서명 (n건): \hrulefill}\par
\else
\begin{itemize}
\item 갱신: 월 1회(동결일 다음 주 PM 회의, 신규 접수 요구와 함께 신규 리스크 등록) + 게이트마다 버전 상향(RACI 26행). 소유자는 자기 리스크의 트리거를 관측하고 월 회의에 보고한다.
\item 잔여 위협 EMV 합 8억 초과 \textbf{예상} 또는 단일 리스크 잔여 EMV 2억 초과 시 5영업일 내 스폰서 예외 보고(헌장 §12 형식). 관리예비비 대상 리스크(R-01·R-04)의 트리거 발동은 즉시 스폰서·프로그램 PMO장 보고.
\item 리스크 실현(이슈 전환) 시 컨틴전시 배정액 집행(⑨ 항목 6 결재), 초과분은 미배정 풀 → 관리예비비 순.
\item 소유자 서명: 13+2건의 소유자가 G1 제출본에 서명한다(Day1 §7 예고). 서명하지 않은 소유자의 리스크는 스폰서에게 소유자 재지정을 요청한다.
\end{itemize}\fi
\end{fieldbox}
\B{\vspace{4mm}\noindent{\footnotesize\begin{tabularx}{\linewidth}{|Y|Y|Y|Y|Y|}\hline
\hd{소유자 서명} & \hd{소유자 서명} & \hd{소유자 서명} & \hd{소유자 서명} & \hd{소유자 서명} \\ \hline
\rd{12mm} & & & & \\ \hline \rd{12mm} & & & & \\ \hline \rd{12mm} & & & & \\ \hline
\end{tabularx}}}{}
\end{document}
